\documentclass[letterpaper]{article} % DO NOT CHANGE THIS
\usepackage[submission]{aaai2027}  % DO NOT CHANGE THIS
% The serif, sans-serif, and monospaced fonts are loaded automatically by
% aaai2027.sty (newtxtext, helvet, courier). DO NOT add \usepackage{times},
% \usepackage{helvet}, \usepackage{courier}, or any other font package.
\usepackage[hyphens]{url}  % DO NOT CHANGE THIS
\usepackage{graphicx} % DO NOT CHANGE THIS
\urlstyle{rm} % DO NOT CHANGE THIS
\def\UrlFont{\rm}  % DO NOT CHANGE THIS
\usepackage{natbib}  % DO NOT CHANGE THIS AND DO NOT ADD ANY OPTIONS TO IT
\usepackage{caption} % DO NOT CHANGE THIS AND DO NOT ADD ANY OPTIONS TO IT
\frenchspacing  % DO NOT CHANGE THIS
\usepackage{amsmath}
\usepackage{amssymb}
\usepackage{booktabs}
\usepackage{multirow}
\usepackage{adjustbox}
\setlength{\textfloatsep}{8pt plus 2pt minus 2pt}
\setlength{\dbltextfloatsep}{8pt plus 2pt minus 2pt}
%
% Keep the \pdfinfo as shown here. There's no need
% for you to add the /Title and /Author tags.
\pdfinfo{
/TemplateVersion (2027.1)
}

% DISALLOWED PACKAGES
% \usepackage{authblk} -- This package is specifically forbidden
% \usepackage{balance} -- This package is specifically forbidden
% \usepackage{CJK} -- This package is specifically forbidden
% \usepackage{float} -- This package is specifically forbidden
% \usepackage{flushend} -- This package is specifically forbidden
% \usepackage{fullpage} -- This package is specifically forbidden
% \usepackage{geometry} -- This package is specifically forbidden
% \usepackage{hyperref} -- This package is specifically forbidden
% \usepackage{navigator} -- This package is specifically forbidden
% (or any other package that embeds links such as navigator or hyperref)
% \usepackage{indentfirst} -- This package is specifically forbidden
% \usepackage{layout} -- This package is specifically forbidden
% \usepackage{multicol} -- This package is specifically forbidden
% \usepackage{nameref} -- This package is specifically forbidden
% \usepackage{savetrees} -- This package is specifically forbidden
% \usepackage{setspace} -- This package is specifically forbidden
% \usepackage{stfloats} -- This package is specifically forbidden
% \usepackage{tabu} -- This package is specifically forbidden
% \usepackage{titlesec} -- This package is specifically forbidden
% \usepackage{tocbibind} -- This package is specifically forbidden
% \usepackage{ulem} -- This package is specifically forbidden
% \usepackage{wrapfig} -- This package is specifically forbidden
% DISALLOWED COMMANDS
% \nocopyright -- Your paper will not be published if you use this command
% \addtolength -- This command may not be used
% \balance -- This command may not be used
% \baselinestretch -- Your paper will not be published if you use this command
% \clearpage -- No page breaks of any kind may be used for the final version of your paper
% \columnsep -- This command may not be used
% \newpage -- No page breaks of any kind may be used for the final version of your paper
% \pagebreak -- No page breaks of any kind may be used for the final version of your paper
% \pagestyle -- This command may not be used
% \tiny -- This is not an acceptable font size.
% \vspace{- -- No negative value may be used in proximity of a caption, figure, table, section, subsection, subsubsection, or reference
% \vskip{- -- No negative value may be used to alter spacing above or below a caption, figure, table, section, subsection, subsubsection, or reference

\newcommand{\method}{LoGTA}

\setcounter{secnumdepth}{0} %May be changed to 1 or 2 if section numbers are desired.

% The file aaai2027.sty is the style file for AAAI Press
% proceedings, working notes, and technical reports.
%

% Title

% Your title must be in mixed case, not sentence case.
% That means all verbs (including short verbs like be, is, using,and go),
% nouns, adverbs, adjectives should be capitalized, including both words in hyphenated terms, while
% articles, conjunctions, and prepositions are lower case unless they
% directly follow a colon or long dash
\title{LoGTA: Uncertainty-Aware Local-to-Global Temporal Semantics for Multimodal Time-Series Anomaly Detection}
% LoGTA: Local-to-Global Temporal Semantics with Uncertainty-Aware Alignment for Multimodal Time-Series Anomaly Detection

\author{
    Anonymous Submission
}
\affiliations{
}

\begin{document}

\maketitle

\begin{abstract}
Multimodal time-series anomaly detection (MTAD) aims to identify abnormal temporal patterns by exploiting complementary information from heterogeneous modalities. 
Recent text-enhanced methods introduce semantic information into time-series modeling, but existing patch-level textual modeling methods derive semantics from isolated local observations, overlooking that similar local variations may occur under different global temporal structures.  
Moreover, similar textual semantics may correspond to diverse temporal patterns, introducing uncertainty into cross-modal correspondence. 
To address these challenges, we propose \textbf{LoGTA}, an uncertainty-aware local-to-global temporal semantic framework for MTAD. 
LoGTA first constructs local behavioral semantics for each temporal patch from its statistical characteristics. 
In parallel, it decomposes the input sequence into trend, seasonal, and residual components to capture global temporal structures. Subsequently, the local patch semantics are enriched with global temporal information, producing more comprehensive semantic representations. LoGTA further estimates the distribution for each temporal patch feature from its enhanced semantic representation, thereby capturing uncertainty in their cross-modal correspondence. The aligned semantic representations then guide masked time-series reconstruction, and reconstruction errors are used to produce anomaly scores. Extensive experiments on six real-world datasets demonstrate that LoGTA achieves competitive performance, while comprehensive ablation studies validate the effectiveness of its key components.
\end{abstract}

% Uncomment the following to link to your code, datasets, an extended version or similar.
% You must keep this block between (not within) the abstract and the main body of the paper.
% Make sure that you do not de-anonymize yourself with these links.
% \begin{links}
%     \link{Code}{https://aaai.org/example/code}
%     \link{Datasets}{https://aaai.org/example/datasets}
%     \link{Extended version}{https://aaai.org/example/extended-version}
% \end{links}
\section{Introduction}
Multimodal Time-series Anomaly Detection (MTAD) aims to identify abnormal patterns by jointly exploiting heterogeneous information associated with time series. It plays an important role in ensuring the reliable operation of industrial equipment, online services, and other continuously monitored systems \cite{su2019robust}. Since anomalous events are rare and difficult to label, existing methods typically learn normal temporal patterns from unlabeled multimodal data and detect anomalies through deviations \cite{audibert2020usad}. However, effectively integrating heterogeneous modalities remains challenging, as different modalities characterize temporal behaviors from distinct perspectives. 
\begin{figure}[t]
    \centering
    \includegraphics[width=\columnwidth]{Figures/motivation.pdf}
    \caption{Similar local patterns may correspond to distinct semantic interpretations under different global contexts.}
    \label{fig:motivation}
\end{figure}

Large Language Models (LLMs) have provided new opportunities to bridge time series and textual modalities. Recent studies have explored different strategies for incorporating textual information into time-series representations. LLMTime~\cite{gruver2023large} transforms time series into a format that can be processed by LLMs. PromptCast~\cite{xue2023promptcast} reformulates time-series forecasting as a natural language generation task. TEST~\cite{sun2024test} aligns time series with textual prototypes to establish semantic correspondence between two modalities. Following these advances, some studies further investigate patch-level
text--temporal modeling. Time-LLM~\cite{jin2024time} reprograms temporal patches into the input space of frozen LLMs for time-series forecasting, while MindTS~\cite{hu2026towards} generates patch-level textual descriptions from temporal statistics via LLMs to provide semantic guidance for time-series reconstruction. Nevertheless, existing patch-level textual modeling methods mainly characterize local temporal behaviors from individual patches, overlooking whether a local variation is consistent with the overall temporal structure.

Indeed, similar local variations in time series may have distinct semantic interpretations under different temporal contexts. As shown in Figure~\ref{fig:motivation}, the selected patches from two time series exhibit similar sharp increases and therefore yield the same local textual description. However, the increase in Series A is consistent with the overall upward trend and is regarded as normal, whereas that in Series B deviates from a stable baseline and is considered anomalous. This indicates that local textual semantics should be interpreted together with the global temporal structure to obtain more complete patch-level semantic representations. Moreover, textual semantics do not uniquely determine temporal representations, as similar semantic descriptions may correspond to diverse temporal patterns. Consequently, deterministic pointwise alignment is overly restrictive and fails to capture the underlying cross-modal uncertainty.


To address these challenges, we propose \textbf{LoGTA}, a \textbf{Lo}cal-to-\textbf{G}lobal \textbf{T}emporal semantic modeling framework with uncertainty-\textbf{A}ware cross-modal alignment for MTAD. Specifically, LoGTA first constructs local behavioral semantics for each patch from its statistical characteristics, while capturing the global temporal structure of the input sequence by decomposing it into trend, seasonal, and residual components. The local patch semantics then interact with the global structural representations to obtain context-aware semantic features. To model the uncertain correspondence between semantic and temporal representations, LoGTA estimates a conditional Gaussian distribution for each temporal patch feature from its enhanced semantic representation. The aligned representations further guide masked time-series reconstruction, and the reconstruction error is used to produce timestamp-level anomaly scores during inference. Our main contributions are summarized as follows:
\begin{itemize}
    \item We propose a local-to-global temporal semantic framework that combines patch-wise behavioral semantics with global trend, seasonal, and residual structures, thereby enriching local semantics with global temporal information.
    
    \item We introduce an uncertainty-aware patch-wise cross-modal alignment mechanism to model the ambiguous correspondence between semantic and temporal representations, thereby capturing the variability of temporal features associated with similar semantics.
    
    \item Extensive experiments on six real-world datasets demonstrate that LoGTA achieves competitive performance. Notably, comprehensive ablation studies further validate the effectiveness of its key components. 
\end{itemize}



\section{Related Work}

\subsubsection{Time-Series Anomaly Detection.}

Time-series anomaly detection (TSAD) assigns anomaly scores to individual timestamps or temporal segments \cite{schmidl2022anomaly}. NPSR \cite{lai2023nominality} combines point-wise and sequential reconstruction, whereas PatchAD \cite{zhong2025patchad} organizes observations into multiscale patches for temporal modeling. Deep TSAD methods are commonly categorized as forecasting-based, reconstruction-based, representation-based, or hybrid approaches, which detect anomalies through prediction errors, reconstruction errors or probabilities, latent-space discrepancies, or combined criteria \cite{zamanzadeh2024deep}. Among reconstruction-based approaches, OmniAnomaly \cite{su2019robust} models multivariate temporal dependencies with stochastic latent variables and estimates reconstruction probabilities, whereas Anomaly Transformer \cite{xu2021anomaly} combines reconstruction error with association discrepancy. Recent studies further enrich numerical modeling with structural information: GCAD \cite{liu2025gcad} captures changes in dynamic Granger-causal relations, SDA2D \cite{sun2026state} models derivatives of latent system states, PatchAD \cite{zhong2025patchad} learns intra- and inter-patch relations, and RobustTAD \cite{gao2020robusttad} integrates seasonal--trend decomposition with a convolutional encoder--decoder. Despite these advances, the relation between local behavior and its broader temporal context remains implicit in numerical residuals, making similar local patterns under different global dynamics difficult to distinguish.

\subsubsection{Multimodal Time-Series Analysis.}
Multimodal time-series analysis integrates temporal data with complementary modalities through modality transference, representation alignment, and cross-modal fusion \cite{jiang2025multi}. VisionTS \cite{chen2025visionts} reformulates time-series forecasting as image reconstruction to leverage a visual masked autoencoder pretrained on natural images, while Time-VLM \cite{zhong2025time} encodes time series as images and integrates temporal, visual, and textual representations through pretrained vision-language models. In text--temporal modeling, PromptCast \cite{xue2023promptcast} and LLMTime \cite{gruver2023large} formulate forecasting as language generation by converting numerical sequences into textual prompts or digit strings. TEST \cite{sun2024test} and Time-LLM \cite{jin2024time} instead adapt temporal representations to frozen language-model embedding spaces through contrastive alignment and prototype-based reprogramming. TimeCMA \cite{liu2025timecma} and M3Time \cite{jia2026m3time} further strengthen text--temporal interaction through cross-modal alignment or the integration of LLM-derived semantic priors with fine-grained temporal features. These methods primarily address forecasting, classification, or representation learning. SigLLM \cite{alnegheimish2024large} and LLMAD \cite{liu2025large} employ textualized time series with prompting or in-context reasoning. MindTS \cite{hu2026towards} more directly combines patch-derived endogenous text with exogenous background information, contrastively aligning temporal and textual patch representations. 
% Despite these advances, existing methods typically rely on deterministic
% patch-level alignment, without
% explicitly accounting for modality-specific discrepancies and the ambiguity
% of cross-modal correspondence. Consequently, uncertainty-aware
% patch-wise cross-modal alignment remains underexplored.
Despite these advances, patch-wise cross-modal alignment that explicitly
accounts for modality-specific discrepancies and ambiguous
correspondence remains underexplored.




\section{Methodology} 
\begin{figure*}[!t]
\centering
\includegraphics[width=\textwidth]{Figures/aaaimi.pdf}
\caption{Overall framework of the proposed \method{}.}
\label{fig:architecture}
\end{figure*}

\subsection{Problem Definition}
Let $\mathbf{X}=[\mathbf{x}_1,\mathbf{x}_2,\ldots,\mathbf{x}_L]\in\mathbb{R}^{L\times C}$ denote a multivariate series, where $L$ is the sequence length, $C$ is the number of variables, and $\mathbf{x}_i\in\mathbb{R}^{C}$ is the observation at timestamp $t_i$. Each observation $\mathbf{x}_i$ corresponds to an anomaly label $y_i\in\{0,1\}$, $y_i=0$ indicates a normal observation and $y_i=1$ indicates an anomaly. These labels are used only for evaluation. In our MTAD setting, textual semantics are derived from $\mathbf{X}$ and combined with its numerical representations to determine whether each observation is normal or anomalous.
%                    Input Time Series X
%                            |
%           +----------------+----------------+
%           |                                 |
%           v                                 v
% Patch Embedding and                Full-window Decomposition
%  Temporal Encoder               Trend / Seasonal / Residual
%           |                                 |
%           |                       Component-specific Prompts
%           |                                 |
%           |                       Prompt-conditioned Encoding
%           |                                 |
%           |                       Structural Context V
%           |                                 |
%           |         Patch Statistics        |
%           |                |                 |
%           |         Endogenous Text          |
%           |                |                 |
%           |          Frozen LLM              |
%           |                |                 |
%           |       Local Semantics U ---------+
%           |                |
%           |      Structure-conditioned
%           |       Semantic Enhancement
%           |                |
%           |               Hc
%           |                |
%           +------ Probabilistic Alignment
%           |                |
%          Hs                |
%           |                |
%           +---- Semantic-conditioned Reconstruction
%                            |
%                       Reconstructed X
%                            |
%                    Reconstruction Error
%                            |
%                      Anomaly Score
\subsection{Framework Overview}
The overall framework of the proposed \textbf{LoGTA} is illustrated in Figure~\ref{fig:architecture}. Given an  multivariate time series $\mathbf{X}$, we first divide it into patch-wise segments $\mathbf{P}$ following MindTS~\cite{hu2026towards}. Local statistical texts are constructed for each patch and encoded by a frozen LLM to obtain patch-level semantic representations. Meanwhile, the complete sequence is decomposed into trend, seasonal, and residual components following STL~\cite{cleveland1990stl}, 
and the component-aware representations are fused to capture global temporal dynamics. The global structural representation is then used to enhance local semantic representations through cross-attention, producing structure-aware semantic representations. Subsequently, a patch-wise semantic distribution alignment module is introduced to probabilistically align semantic and temporal representations in a shared latent space. Finally, the aligned semantic representations guide masked time-series reconstruction, and the reconstruction error is used as the anomaly score. 
% We omit  the batch dimension and present latent dimensions for one variable; the same operations are applied independently to all $C$ variables.
% The proposed framework reconstructs each input window by jointly modeling numerical temporal patterns and structured textual descriptions derived from the time series itself. As illustrated in Figure~\ref{fig:architecture}, \method{} consists of three main stages: (1) temporal patch encoding with masking, (2) dual-view textual modality construction, and (3) text-temporal alignment for masked reconstruction.

% Given an input window $\mathbf{X}\in \mathbb{R}^{L\times C}$, we first apply window-level normalization:
% \begin{equation}
%     \bar{\mathbf{X}} = \frac{\mathbf{X}-\mathbf{m}}{\boldsymbol{\sigma}}, \quad \boldsymbol{\sigma} = \sqrt{\mathrm{Var}(\mathbf{X})+\epsilon},
% \end{equation}
% where $\mathbf{m}$ and $\boldsymbol{\sigma}$ are computed along the temporal dimension and $\epsilon$ is a small constant for numerical stability. The normalized sequence is divided into patches with size $p$ and stride $r$, yielding $N_p = \lfloor(L-p)/r\rfloor + 1$ patches. Each patch is projected to dimension $d$ via a learnable linear projection $\mathbf{W}_{\mathrm{patch}} \in \mathbb{R}^{p \times d}$, followed by additive positional encoding. The resulting patch representation is $\mathbf{P}\in \mathbb{R}^{BC\times N_p\times d}$, where $B$ is the batch size and $d$ is the hidden dimension. A random patch mask is applied to construct incomplete temporal evidence. Both the complete and masked patch sequences are encoded by a shared Transformer encoder $E_t$, where $\mathrm{Mask}(\cdot)$ randomly zeros out a fraction of patch tokens:
% \begin{equation}
%     \mathbf{H}^{t}=E_t(\mathbf{P}), \quad \widetilde{\mathbf{H}}^{t} = E_t(\mathrm{Mask}(\mathbf{P})),
% \end{equation}
% where $\mathbf{H}^{t}$ serves as the complete temporal representation for alignment and $\widetilde{\mathbf{H}}^{t}$ serves as the masked temporal representation for reconstruction.

\subsection{Local-to-Global Semantic Representation}
We first construct patch-level behavioral texts from local statistical characteristics and lag dependence, and encode them as local semantic representations. We then decompose the complete input sequence into trend, seasonal, and residual components; these structural representations are incorporated into the local semantics through cross-attention, yielding globally informed patch-level semantic representations.

\noindent\textbf{Local Statistical Text Generation.}
Given the patch-wise sequence
$\mathbf{P}=[\mathbf{p}_1,\mathbf{p}_2,\ldots,\mathbf{p}_N]$ for time series $\mathbf{x}$, each patch satisfies $\mathbf{p}_i\in\mathbb{R}^{p \times C}(i=1,2,\cdots, N)$ and $N=\lfloor(L-p)/\delta\rfloor+1$ for the patch size $p$ and stride $\delta$. 
% A patch embedder $E_p$ maps each patch to a $d$-dimensional
% token. A shared temporal encoder $E_t$ processes the original and masked
% patch tokens, yielding
% $\mathbf{H}^{o},\widetilde{\mathbf{H}}^{\gamma}\in\mathbb{R}^{N\times d}$:
% \begin{equation}
%     \mathbf{H}^{o}=E_t(E_p(\mathbf{P})),\quad
%     \widetilde{\mathbf{H}}^{\gamma}
%     =E_t\left(\operatorname{Mask}_{\gamma}(E_p(\mathbf{P}))\right),
% \end{equation}
% where $\operatorname{Mask}_{\gamma}(\cdot)$ masks a fraction $\gamma$ of
% the patch tokens and the superscript $o$ denotes the original unmasked
% input. 
For each temporal patch $\mathbf{p}_i$, we summarize its local behavior using several statistical descriptors, including the minimum, maximum, and median values. We also estimate its dominant lags through autocorrelation and derive a coarse trend descriptor from the first-order differences within the patch. These descriptors are combined with the dataset and task descriptions and verbalized as a patch-level natural language prompt $c_i$. Accordingly, the prompts for all $N$ patches are denoted by $\mathcal{C}_{\mathrm{text}}=[c_1,c_2,\ldots,c_N]$. After being padded or truncated to a maximum length of $M$ tokens, these prompts are independently encoded by a frozen Large Language Model (LLM):
\begin{equation}
    \widetilde{\mathbf{H}}^{\mathrm{llm}}
    =\mathrm{LLM}(\mathcal{C}_{\mathrm{text}})
    \in
    \mathbb{R}^{N\times M\times d_{\mathrm{llm}}},
\end{equation}
where $M$ denotes the maximum prompt length and $d_{\mathrm{llm}}$ is the
hidden dimension of the LLM. To obtain a single semantic representation for
each patch, we first project the token-level hidden states into a
$d$-dimensional space:
\begin{equation}
    \mathbf{H}^{\mathrm{llm}}
    =
    \widetilde{\mathbf{H}}^{\mathrm{llm}}\mathbf{W}_{\mathrm{llm}}
    \in \mathbb{R}^{N\times M\times d},
\end{equation}
where
$\mathbf{W}_{\mathrm{llm}}\in
\mathbb{R}^{d_{\mathrm{llm}}\times d}$.
The projected token representations associated with each patch are then
concatenated and mapped into a patch-level semantic feature:
\begin{equation}
    \mathbf{U}^{\mathrm{loc}}
    =
    \left(
    \operatorname{Flatten}_{\mathrm{token}}
    \left(
        \mathbf{H}^{\mathrm{llm}}
    \right)
    \right)
    \mathbf{W}_{\mathrm{u}}
    \in\mathbb{R}^{N\times d},
\end{equation}
where $\operatorname{Flatten}_{\mathrm{token}}(\cdot)$ preserves the patch
dimension while flattening the $M$ token representations of each prompt,  $\mathbf{W}_{\mathrm{u}}\in\mathbb{R}^{Md\times d}$, and
$\mathbf{U}^{\mathrm{loc}}$ denotes the resulting local semantic
representations of all temporal patches.


% Here
% 设输入序列为 $\mathbf{X}\in\mathbb{R}^{B\times L\times C}$。代码
% \texttt{self.avg(x.permute(0, 2, 1)).permute(0, 2, 1)}
% 首先将输入由 $[B,L,C]$ 转换为 $[B,C,L]$，使 \texttt{AvgPool1d} 能够沿时间维 $L$ 执行滑动平均，随后再将结果转换回 $[B,L,C]$。因此，移动平均结果满足
% \[
% \operatorname{MA}(\mathbf{X})\in\mathbb{R}^{B\times L\times C},
% \]
% 忽略批次维度后，可写为
% \[
% \operatorname{MA}(\mathbf{X})\in\mathbb{R}^{L\times C}.
% \]
% 趋势映射层定义为
% \texttt{self.map\_trend = nn.Linear(configs.seq\_len, configs.seq\_len)}，
% 其输入维度和输出维度均为 $L$，因此权重矩阵和偏置向量分别满足
% \[
% \mathbf{W}_{t}\in\mathbb{R}^{L\times L},
% \qquad
% \mathbf{b}_{t}\in\mathbb{R}^{L}.
% \]
% 在代码
% \texttt{trend = self.map\_trend(trend.transpose(1, 2)).transpose(1, 2)}
% 中，移动平均结果首先由 $[B,L,C]$ 转置为 $[B,C,L]$，随后 \texttt{Linear}$(L,L)$ 对每个通道沿时间维执行线性映射，输出维度仍为 $[B,C,L]$，最后再转置回 $[B,L,C]$。因此，忽略批次维度后，趋势分量可表示为
% \[
% \mathbf{X}^{\mathrm{tr}}
% =
% \mathbf{W}_{t}\operatorname{MA}(\mathbf{X})
% +
% \mathbf{b}_{t}\mathbf{1}_{C}^{\top},
% \]
% 其中 $\mathbf{1}_{C}\in\mathbb{R}^{C}$ 为全一向量，其维度关系为
% \[
% \underbrace{\mathbf{W}_{t}}_{L\times L}
% \underbrace{\operatorname{MA}(\mathbf{X})}_{L\times C}
% +
% \underbrace{\mathbf{b}_{t}}_{L\times 1}
% \underbrace{\mathbf{1}_{C}^{\top}}_{1\times C}
% =
% \underbrace{\mathbf{X}^{\mathrm{tr}}}_{L\times C}.
% \]
% 需要说明的是，公式中不需要再对 $\operatorname{MA}(\mathbf{X})$ 进行转置；代码中的两次转置只是为了将时间维 $L$ 放置在最后一维，使 PyTorch 的 \texttt{nn.Linear} 能够沿该维度完成映射。



\noindent\textbf{Global Structure-aware Series Representation.}
Although the patch-wise semantic representation $\mathbf{U}^{\mathrm{loc}}$ describes local statistical variations, it does not explicitly capture the structural information of the complete input sequence, such as global trend and periodicity. Following the STL decomposition strategy~\cite{cao2024tempo}, we decompose the input sequence as follows:
\begin{equation}
    \mathbf{X}=\mathbf{X}^{t}+\mathbf{X}^{s}+\mathbf{X}^{r},
\end{equation}
where $\mathbf{X}^{t}$, $\mathbf{X}^{s}$, and $\mathbf{X}^{r}$ denote the trend, seasonal, and residual components, respectively. Specifically, the trend component is first extracted by a moving-average
filter and then refined through a learnable temporal projection. Let
$\operatorname{MA}_{k}(\cdot)$ denote a moving-average operator with
kernel size $k$. The trend component is defined as:
\begin{equation}
    \mathbf{X}^{t}=\mathbf{W}_{t}[\operatorname{MA}_{k}(\mathbf{X})] + \mathbf{b}_{t},
\end{equation}
where $\mathbf{W}_{t}\in\mathbb{R}^{L\times L}$ and
$\mathbf{b}_{t}\in\mathbb{R}^{L}$ are learnable parameters.
The seasonal component is then estimated from the detrended sequence
through a two-layer feed-forward network:
\begin{equation}
    \mathbf{X}^{s}
    =
    \mathbf{W}_{s}^2
    \,
    \operatorname{ReLU}
    \left(
    \mathbf{W}_{s}^1
    (    \mathbf{X}
    -
    \mathbf{X}^{t})
    +
    \mathbf{b}_{s}^1
    \right)
    +
    \mathbf{b}_{s}^2
    ,
\end{equation}
where $\mathbf{W}_{s}^1\in\mathbb{R}^{4L\times L}$,
$\mathbf{W}_{s}^2\in\mathbb{R}^{L\times4L}$,
$\mathbf{b}_{s}^1\in\mathbb{R}^{4L}$, and
$\mathbf{b}_{s}^2\in\mathbb{R}^{L}$ are learnable parameters.
Finally, the residual component $\mathbf{X}^{r}$ is obtained by subtracting the trend
and seasonal components from the original sequence.


% \begin{equation}
%     \mathbf{X}^{r}
%     =
%     \mathbf{X}
%     -
%     \mathbf{X}^{t}
%     -
%     \mathbf{X}^{s}.
% \end{equation}

For each component $q\in\{t,s,r\}$ and each variable, we first partition $\mathbf{X}^{q}$ into patches using the same patch size $p$ and stride $\delta$ as the original sequence, and project the resulting patches into $d$-dimensional embeddings $\mathbf{P}^{q}\in\mathbb{R}^{N\times d}$. Then, a component-specific role prompt is constructed for each component, such as "Reconstruct the time series using its trend, seasonal, or residual component." After LLM encoding and feature projection, the
corresponding prompt representation is denoted by
$\mathbf{R}^{q}\in\mathbb{R}^{M\times d}$. The prompt representation and component patch embeddings are
concatenated along the token dimension and processed by a component Transformer encoder $\mathcal{T}_e$ shared across the three components:
\begin{equation}
    \mathbf{V}^{q}
    =
    \Pi_{\mathrm{patch}}
    \left(
    \mathcal{T}_e
    \left(
    [\mathbf{R}^{q}; \mathbf{P}^{q}]
    \right)
    \right),
    \qquad
    q\in\{t,s,r\},
\end{equation}
where $[\cdot;\cdot]$ denotes concatenation along the token dimension, $\Pi_{\mathrm{patch}}\left(\cdot\right)$ removes the prompt tokens and retains the $N$ patch-wise tokens for each component $q$, and $\mathbf{V}^{q}\in\mathbb{R}^{N \times d}$ contains the output representations at the $N$ patch positions. Then the three component representations are concatenated along the feature
dimension and fused through a learnable projection:
\begin{equation}
    \mathbf{V}^{\mathrm{glo}}
    =
    \left[
    \mathbf{V}^{t}
    ;
    \mathbf{V}^{s}
    ;
    \mathbf{V}^{r}
    \right]
    \mathbf{W}_{c},
\end{equation}
where $[\cdot;\cdot]$ denotes feature concatenation,
$\mathbf{W}_{c}\in\mathbb{R}^{3d\times d}$ is a learnable parameter, and
$\mathbf{V}^{\mathrm{glo}}\in\mathbb{R}^{N\times d}$ denotes
the fused global temporal representations of the $N$ patches.

% 在这个地方写一句话表示两个
\noindent\textbf{Global-Guided Local Semantic Enhancement.} The global representation $\mathbf{V}^{\mathrm{glo}}$ enhances the local semantics
$\mathbf{U}^{\mathrm{loc}}$ through cross-attention:
\begin{align}
    \makebox[0pt][l]{\ensuremath{\mathbf{Z}_1}}
    \phantom{\mathbf{H}^{\mathrm{sem}}}
    &=
    \operatorname{LN}\!\left(
    \mathbf{U}^{\mathrm{loc}}
    +\mathcal{A}(\mathbf{U}^{\mathrm{loc}},\mathbf{V}^{\mathrm{glo}},\mathbf{V}^{\mathrm{glo}})
    \right),\\
    \mathbf{H}^{\mathrm{sem}}
    &=
    \operatorname{LN}\!\left(
    \mathbf{Z}_1+\operatorname{FFN}(\mathbf{Z}_1)
    \right).
\end{align}
Here, $\mathcal{A}$ denotes cross-attention with
$\mathbf{U}^{\mathrm{loc}}$ as the query and $\mathbf{V}^{\mathrm{glo}}$ as the key and
value, while $\operatorname{LN}$ and $\operatorname{FFN}$ denote layer
normalization and a feed-forward network, respectively. The resulting representation
$\mathbf{H}^{\mathrm{sem}}\in\mathbb{R}^{N\times d}$
denotes the globally enhanced semantic features of the $N$ temporal patches.
\subsection{Patch-wise Cross-Modal Distribution Alignment}
To account for uncertainty in semantic--temporal relationships, 
we model the correspondence of two cross-modal representations via Gaussian distributions. For patch-wise enhanced semantic features $\mathbf{H}^{\mathrm{sem}}$, we map them into probability space through a linear projection:
% Specifically, the semantic representation predicts the mean and log-variance of the corresponding temporal representation:
\begin{equation}
\boldsymbol{\mu}
=\mathbf{H}^{\mathrm{sem}}\mathbf{W}_{\mu}+\mathbf{b}_{\mu},
\quad
\log\boldsymbol{\sigma}^{2}
=\mathbf{H}^{\mathrm{sem}}\mathbf{W}_{\sigma}+\mathbf{b}_{\sigma},
\label{eq:distribution_parameters}
\end{equation}
where $\mathbf{W}_{\mu}, \mathbf{W}_{\sigma} \in \mathbb{R}^{d \times d}$, $\mathbf{b}_{\mu}, \mathbf{b}_{\sigma} \in \mathbb{R}^{d}$, $\mathbf{H}^{\mathrm{sem}} = [\mathbf{h}^{\mathrm{sem}}_1, \mathbf{h}^{\mathrm{sem}}_2, \cdots, \mathbf{h}^{\mathrm{sem}}_N]$, and $\boldsymbol{\mu}, \boldsymbol{\sigma}^2 \in
\mathbb{R}^{N \times d}$ represent the patch-wise mean and variance of the semantic representation, respectively.
For numerical stability, the predicted variance
$\boldsymbol{\sigma}^{2}$ is confined within a predefined range.

Conditioned on the semantic representation $\mathbf{h}^{\mathrm{sem}}_{i}$ of the $i$-th patch, the corresponding temporal representation $\mathbf{h}_i^o \in \mathbb{R}^d$ from temporal encoder is modeled by the following conditional multivariate Gaussian distribution:
\begin{equation}
    p
    \left(
    \mathbf{h}^{o}_{i}
    \mid
    \mathbf{h}^{\mathrm{sem}}_{i}
    \right)
    =
    \mathcal{N}
    \left(
    \mathbf{h}^{o}_{i};
    \boldsymbol{\mu}_{i},
    \operatorname{diag}
    \left(
    \boldsymbol{\sigma}_{i}^2
    \right)
    \right),
\end{equation}
which encourages the two-modal features to fall within a high-probability region while accounting for cross-modal variability through the predicted variance. The alignment objective for all patches is the negative log-likelihood (NLL):
\begin{equation}
    \mathcal{L}_{\mathrm{align}}
    =
    \frac{1}{2Nd}
    \sum_{i=1}^{N}
    \sum_{j=1}^{d}
    \Big(
    \log\sigma_{i,j}^2
    +\frac{(h^{o}_{i,j}-\mu_{i,j})^2}
    {\sigma_{i,j}^2}
    \Big).
\end{equation}
\subsection{Semantic-guided Time-Series Reconstruction and Anomaly Scoring}
The aligned semantic features
$\mathbf{H}^{\mathrm{sem}}$ are then used to guide the recovery of the masked temporal patches. Specifically, the patch-wise masked temporal representation 
$\widetilde{\mathbf{H}}^{\gamma}$ attends to
$\mathbf{H}^{\mathrm{sem}}$ through the following cross-attention:
\begin{align}
    \makebox[0pt][l]{\ensuremath{\mathbf{Z}_{2}}}
    \phantom{\hat{\mathbf{H}}^{\gamma}}
    &=
    \operatorname{LN}\!\left(
    \widetilde{\mathbf{H}}^{\gamma}
    +\mathcal{A}(
    \widetilde{\mathbf{H}}^{\gamma},
    \mathbf{H}^{\mathrm{sem}},
    \mathbf{H}^{\mathrm{sem}})
    \right),\\
    \hat{\mathbf{H}}^{\gamma}
    &=
    \operatorname{LN}\!\left(
    \mathbf{Z}_{2}+\operatorname{FFN}(\mathbf{Z}_{2})
    \right).
\end{align}
The reconstructed sequence is $\widehat{\mathbf{X}}$ obtained by projecting patch-wise mask temporal representation $\hat{\mathbf{H}}^{\gamma}$ into data space.  The reconstruction loss is then defined as the mean squared error (MSE) between the input and reconstructed sequence:
\begin{equation}
    \mathcal{L}_{\mathrm{rec}}
    =
    \frac{1}{LC}\sum_{i=1}^{L}
    \left\|
    \mathbf{x}_i-\widehat{\mathbf{x}}_i
    \right\|_{2}^{2}.
\end{equation}
During inference, $\mathcal{L}_{\mathrm{rec}}$ serves as the
anomaly score for the sequence; a larger value indicates a higher
likelihood that the window is anomalous.
The overall training objective combines the reconstruction and
distribution-alignment losses:
\begin{equation}
    \mathcal{L}
    =
    \mathcal{L}_{\mathrm{rec}}
    +
    \lambda
    \mathcal{L}_{\mathrm{align}},
\end{equation}
where $\lambda$ controls the contribution of the alignment objective.



\begin{table*}[!t]
\centering
\resizebox{\textwidth}{!}{
\begin{tabular}{l|c|c|c|c|c|c|c|c|c|c|c|c|c|c}
\toprule
\textbf{Dataset} & \textbf{Metric} & \textbf{LoGTA} & \textbf{MindTS} & \textbf{iTrans.} & \textbf{LightTS} & \textbf{NLinear} & \textbf{Linear} & \textbf{TimesNet} & \textbf{FiLM} & \textbf{Informer} & \textbf{Reformer} & \textbf{Trans.} & \textbf{MICN} & \textbf{FEDformer} \\
\midrule
\multirow{3}{*}{Weather} & Aff-F & \textbf{72.33} & 70.79 & \underline{72.25} & 60.19 & 61.44 & 61.01 & 70.51 & 68.21 & 62.28 & 65.03 & 62.99 & 64.22 & 53.70 \\
& V-PR & 48.53 & 47.91 & \textbf{50.61} & 31.29 & 36.87 & 38.86 & 47.56 & 44.26 & \underline{49.16} & 48.84 & 48.43 & 41.89 & 39.63 \\
& V-ROC & 77.07 & 76.41 & \textbf{78.15} & 60.57 & 70.22 & 67.62 & 77.84 & 77.87 & 75.29 & \underline{77.90} & 76.66 & 67.90 & 61.84 \\
\midrule
\multirow{3}{*}{Energy} & Aff-F & \textbf{65.22} & 63.10 & 64.00 & 56.68 & \underline{64.52} & 58.69 & 64.42 & 59.37 & 56.55 & 52.89 & 55.75 & 64.14 & 56.51 \\
& V-PR & \underline{45.83} & 41.10 & 34.94 & 27.82 & 31.76 & 31.55 & \textbf{48.16} & 38.77 & 33.89 & 33.26 & 33.69 & 45.13 & 34.74 \\
& V-ROC & \textbf{71.96} & 67.16 & 62.62 & 50.29 & 58.89 & 55.34 & 69.68 & 60.67 & 57.77 & 54.97 & 56.51 & \underline{71.69} & 57.51 \\
\midrule
\multirow{3}{*}{Environment} & Aff-F & \underline{65.21} & \textbf{66.80} & 61.60 & 57.63 & 59.04 & 55.06 & 61.79 & 63.59 & 50.78 & 57.78 & 54.49 & 57.12 & 50.68 \\
& V-PR & \underline{53.69} & \textbf{54.87} & 39.47 & 14.76 & 25.03 & 14.63 & 27.17 & 36.37 & 13.74 & 15.19 & 8.46 & 11.83 & 7.20 \\
& V-ROC & \underline{91.81} & \textbf{93.13} & 88.13 & 63.42 & 78.45 & 61.95 & 73.57 & 84.48 & 62.57 & 62.85 & 46.92 & 58.87 & 39.84 \\
\midrule
\multirow{3}{*}{KR} & Aff-F & \textbf{79.31} & 78.03 & 74.86 & 63.16 & 69.29 & 67.37 & \underline{79.29} & 75.92 & 70.94 & 69.44 & 62.99 & 59.32 & 69.72 \\
& V-PR & \textbf{55.53} & \underline{54.82} & 46.22 & 27.94 & 31.04 & 33.17 & 54.62 & 44.52 & 50.08 & 39.60 & 31.13 & 19.82 & 43.74 \\
& V-ROC & \textbf{93.29} & \underline{91.56} & 86.85 & 69.58 & 75.21 & 78.22 & 91.26 & 84.71 & 87.06 & 77.76 & 62.36 & 49.74 & 79.92 \\
\midrule
\multirow{3}{*}{EWJ} & Aff-F & \textbf{72.67} & 70.43 & 64.84 & 58.69 & 60.66 & 62.29 & \underline{71.73} & 57.60 & 69.22 & 60.50 & 55.90 & 47.66 & 64.73 \\
& V-PR & \textbf{48.52} & 43.61 & 29.07 & 22.21 & 32.51 & 25.91 & \underline{46.66} & 20.08 & 32.36 & 25.36 & 19.06 & 16.35 & 21.28 \\
& V-ROC & \textbf{82.70} & \underline{80.44} & 72.06 & 61.54 & 73.53 & 66.25 & 80.04 & 60.93 & 72.58 & 63.14 & 54.49 & 52.22 & 58.61 \\
\midrule
\multirow{3}{*}{MDT} & Aff-F & \textbf{69.25} & \underline{69.17} & 64.60 & 55.29 & 57.30 & 54.12 & 67.26 & 66.67 & 62.81 & 62.66 & 53.99 & 51.27 & 60.63 \\
& V-PR & \textbf{59.14} & \underline{57.62} & 39.35 & 22.15 & 27.08 & 20.20 & 50.66 & 49.76 & 33.62 & 33.55 & 17.20 & 17.44 & 27.51 \\
& V-ROC & \textbf{81.65} & \underline{80.99} & 69.53 & 54.65 & 60.81 & 49.33 & 75.99 & 74.44 & 63.82 & 62.32 & 41.11 & 43.75 & 59.33 \\
\bottomrule
\end{tabular}
}
\caption{Performance comparison of \method{} with unimodal and LLM-based baselines on six real-world datasets. Aff-F is computed using the mean-threshold protocol, whereas V-PR and V-ROC are threshold-independent. All results are reported as percentages, with the best and second-best results highlighted in bold and underlined, respectively.}
\label{tab:main_results}
\end{table*}


% 1. 局部-全局结构感知语义建模：融合局部统计文本和全局结构语义，提升patch级语义表示的完整性，使其同时感知局部变化与整体动态模式。
% 2. 概率化跨模态语义分布对齐：建模文本语义到时间表示之间的条件分布关系，显式刻画跨模态对应的不确定性，避免确定性特征匹配带来的过度约束。





\section{Experiments}



\subsection{Experimental Setup}

\noindent\textbf{Datasets.}
We evaluate our model on six MTAD datasets used in MindTS~\cite{hu2026towards}: Weather, Energy, Environment, KR, EWJ, and MDT. Energy and Environment are adapted from Time-MMD~\cite{liu2024time}, while KR, EWJ, and MDT are constructed from the financial data provided by FNSPID~\cite{dong2024fnspid}. Weather is a temperature--humidity dataset released with MindTS. In terms of the numerical time-series modality, Environment, KR, EWJ, and MDT are univariate, whereas Energy and Weather contain 9 and 4 variables, respectively.








\noindent\textbf{Baselines.}
We compare \method{} with Linear and NLinear~\cite{zeng2023transformers}; TimesNet~\cite{wu2022timesnet}, MICN~\cite{wang2023micn}, FiLM~\cite{zhou2022film}, and LightTS~\cite{zhang2022morefastmultivariatetime}; five Transformer-based models: Transformer~\cite{vaswani2017attention}, iTransformer~\cite{liu2024itransformer}, Informer~\cite{zhou2021informer}, Reformer~\cite{kitaev2020reformer}, and FEDformer~\cite{zhou2022fedformer}; and the multimodal model MindTS~\cite{hu2026towards}.





\noindent\textbf{Evaluation Metrics.}
We report affiliation-based F1 (Aff-F)~\cite{huet2022local} under the benchmark's mean-threshold protocol. VUS-PR and VUS-ROC~\cite{paparrizos2022volume}, abbreviated V-PR and V-ROC, provide threshold-independent, range-aware evaluation.

\noindent\textbf{Implementation Details.}
\method{} is implemented in PyTorch with a frozen Qwen3-1.7B backbone and trained for five epochs using Adam with an initial learning rate of $10^{-4}$. The alignment weight of the Gaussian NLL objective is set to $\lambda=1$. Detailed dataset-specific configurations and the computing environment
are provided in the Appendix.

\subsection{Overall Performance}

\noindent To evaluate the overall anomaly detection performance of \method{}
across datasets with different temporal characteristics, we compare it with
four groups of representative baselines. Linear and NLinear serve as linear
baselines; TimesNet, MICN, FiLM, and LightTS model multi-periodicity,
multi-scale contexts, frequency-enhanced historical representations, and
sampling-based temporal dependencies, respectively; Transformer, Informer,
Reformer, FEDformer, and iTransformer employ different attention-based
architectures to model temporal, frequency-domain, or inter-variable
dependencies; and MindTS is a text-enhanced multimodal method. As reported in
Table~\ref{tab:main_results}, \method{} achieves competitive Aff-F results
across the six benchmarks. Compared with the strongest linear baseline,
\method{} yields relative improvements of 20.9\% on MDT and 16.7\% on EWJ,
highlighting the limitations of linear mappings in characterizing complex
patch behavior. Against TimesNet, MICN, FiLM, and LightTS, \method{} improves
over the strongest method by 3.0\% on MDT and 2.6\% on Weather. These gains
are consistent with local-to-global structure-aware semantic modeling, providing
explicit descriptions of local variations together with their overall
temporal context, thereby complementing specialized numerical structures.
Compared with the strongest Transformer-based baseline, the relative gains
reach 7.2\% on MDT, 5.9\% on KR, and 5.0\% on EWJ, suggesting that explicit
patch-level semantics provide complementary information beyond
attention-based numerical dependency modeling. Compared with MindTS,
\method{} improves Aff-F by 3.4\% on Energy and 3.2\% on EWJ. This advantage
may benefit from deriving both local and global semantic information from the
input series and modeling the conditional distribution of temporal
representations given semantic representations. In particular, the estimated variance explicitly captures uncertainty in patch-wise cross-modal correspondences, avoiding overly strong constraints on uncertain pairs.


\begin{table*}[!t]
\centering
\resizebox{\textwidth}{!}{
\begin{tabular}{l|ccc|ccc|ccc|ccc|ccc|ccc}
\toprule
\multirow{2}{*}{\textbf{Model}} & \multicolumn{3}{c}{\textbf{Weather}} & \multicolumn{3}{c}{\textbf{Energy}} & \multicolumn{3}{c}{\textbf{Env.}} & \multicolumn{3}{c}{\textbf{KR}} & \multicolumn{3}{c}{\textbf{EWJ}} & \multicolumn{3}{c}{\textbf{MDT}} \\
\cmidrule(lr){2-4}\cmidrule(lr){5-7}\cmidrule(lr){8-10}\cmidrule(lr){11-13}\cmidrule(lr){14-16}\cmidrule(lr){17-19}
& Aff-F & V-PR & V-ROC & Aff-F & V-PR & V-ROC & Aff-F & V-PR & V-ROC & Aff-F & V-PR & V-ROC & Aff-F & V-PR & V-ROC & Aff-F & V-PR & V-ROC \\
\midrule
Qwen3-1.7B & \textbf{72.33} & \textbf{48.53} & \underline{77.07} & 65.22 & 45.83 & 71.96 & \textbf{65.21} & \textbf{53.69} & \textbf{91.81} & \underline{79.31} & 55.53 & 93.29 & \underline{72.67} & 48.52 & \underline{82.70} & 69.25 & 59.14 & 81.65 \\
Qwen2.5-1.5B & 72.05 & 47.52 & 76.65 & 64.56 & 46.40 & 73.21 & 63.13 & 48.35 & \underline{87.36} & 78.94 & \underline{58.47} & 93.41 & 72.31 & \underline{50.83} & 82.37 & 69.67 & 52.35 & 78.12 \\
SmolLM2-1.7B & \underline{72.13} & 47.29 & 76.69 & 63.74 & 43.64 & 70.02 & 64.73 & 50.46 & 86.26 & \textbf{80.14} & \textbf{58.72} & \textbf{94.57} & 72.54 & 45.80 & 82.24 & \textbf{69.87} & \underline{60.35} & \underline{82.22} \\
TinyLlama-1.1B & 71.65 & 47.84 & 76.64 & \underline{65.60} & \textbf{49.23} & \textbf{74.91} & \underline{65.05} & 49.67 & 85.88 & 78.67 & 55.83 & 93.40 & 72.12 & 44.21 & 80.88 & 69.48 & 58.93 & 81.50 \\
Phi-1.5 & 71.54 & 47.77 & 77.04 & 65.31 & 47.08 & 73.32 & 64.53 & 50.38 & 86.59 & 78.74 & 53.62 & 93.44 & 71.13 & 46.49 & 81.84 & 69.63 & 58.31 & 80.40 \\
OLMo-1B & 71.35 & 46.67 & 76.58 & 65.17 & \underline{47.32} & \underline{74.61} & 64.07 & 50.47 & 86.73 & 77.06 & 46.79 & 89.98 & \textbf{73.06} & \textbf{51.35} & \textbf{84.69} & \underline{69.76} & \textbf{60.95} & \textbf{82.84} \\
MobileLLaMA-1.4B & 71.82 & \underline{48.27} & \textbf{77.33} & \textbf{65.63} & 45.43 & 72.77 & 64.54 & \underline{51.08} & 86.57 & 78.32 & 52.44 & \underline{93.56} & 70.55 & 44.52 & 80.91 & 69.57 & 59.32 & 81.33 \\
\bottomrule
\end{tabular}
}
\caption{Comparison of different frozen LLM backbones. Aff-F is computed using the mean-threshold protocol, whereas V-PR and V-ROC are threshold-independent. }
\label{tab:llm_comparison}
\end{table*}



\begin{figure}[!t]
\centering
\includegraphics[width=\linewidth]{Figures/ablation_aff_lollipop_nature.pdf}
\caption{Aff-F ablation results on four datasets.
(a) Component ablations: \textit{w/o Align.} removes Gaussian NLL alignment;
\textit{w/o Sem.} removes the entire local-to-global semantic branch;
\textit{w/o Local} removes local statistical semantics and their
global enhancement; and \textit{w/o Global} removes global
structural representations and their guidance to local semantics.
(b) Comparison of Gaussian NLL, cosine similarity, and MSE alignment losses.
Env. denotes Environment.}
\label{fig:ablation_aff}
\end{figure}


\begin{figure}[!t]
\centering
\includegraphics[width=1.0\columnwidth]{Figures/hyper.pdf}
\par\smallskip
\makebox[0.5\linewidth][c]{\textbf{(a)} Mask Ratio $\gamma$}%
\makebox[0.5\linewidth][c]{\textbf{(b)} Alignment Weight $\lambda$}
\caption{Sensitivity analysis of the mask ratio $\gamma$ and the alignment weight $\lambda$. Dashed lines indicate default settings.}
\label{fig:hyperparameter}
\end{figure}

\begin{figure}[!t]
\centering
\includegraphics[width=1\linewidth]{Figures/noise_robustness_lollipop.pdf}
\caption{Model performance under different levels of noise.}
\label{fig:test_noise_robustness}
\end{figure}


\begin{figure*}[!t]
\centering
\includegraphics[width=\textwidth]{Figures/detection_case.pdf}
\caption{Qualitative anomaly detection results on six datasets. The top row shows the input time series, with shaded regions indicating ground-truth anomaly intervals. The bottom row presents the corresponding anomaly scores.}
\label{fig:anomaly_score_grid}
\end{figure*}




\begin{figure}[!t]
\centering
\includegraphics[width=\linewidth]{Figures/text_to_time.pdf}
\caption{Gaussian NLL alignment analysis on EWJ. 
(a) Alignment NLL and predicted uncertainty over time. 
(b) NLL distributions for normal and anomalous patches.}
\label{fig:gaussian_nll}
\end{figure}


\subsection{Ablation Study}

\noindent To assess the contributions of local-to-global structure-aware semantic modeling and patch-wise probabilistic cross-modal alignment, we conduct component and alignment-objective ablations. As shown in Figure~\ref{fig:ablation_aff}(a), \textit{w/o Sem.} causes the largest and most consistent performance degradation, indicating that semantic information provides an important complement to numerical reconstruction. \textit{w/o Local} generally produces a larger decrease than \textit{w/o Global}, suggesting that local statistics and lag dependencies provide the primary patch-level semantics, while global information incorporates the overall temporal dynamics to produce a more complete representation of each patch. The lower performance of \textit{w/o Align.} shows that semantic construction alone is insufficient and that explicitly modeling the conditional distribution of temporal representations given semantic representations contributes to cross-modal correspondence. Figure~\ref{fig:ablation_aff}(b) compares different alignment objectives. Cosine similarity and MSE perform deterministic feature matching, whereas Gaussian NLL alignment models each patch-wise correspondence using a mean and a variance, allowing the alignment penalty to account for correspondence uncertainty and avoid imposing overly strong constraints on uncertain patch pairs.

\subsection{LLM Backbone Analysis}

\noindent To examine whether \method{} depends on a particular LLM, we replace Qwen3-1.7B with six frozen LLM backbones of comparable scale while keeping the remaining settings unchanged. As shown in Table~\ref{tab:llm_comparison}, all backbones yield comparable performance across the six datasets, although their relative advantages vary by dataset and metric. Qwen3-1.7B performs well on Weather and Environment, whereas SmolLM2-1.7B and OLMo-1B obtain stronger results on KR and EWJ, respectively. The absence of substantial degradation after backbone replacement indicates that the proposed semantic modeling and probabilistic alignment are not tied to a specific LLM architecture or model family.




\subsection{Hyperparameter Sensitivity Analysis}

\noindent To verify whether local-to-global semantic modeling can reliably support masked reconstruction and whether probabilistic cross-modal alignment depends on a carefully tuned loss weight, we analyze the effects of the mask ratio $\gamma$ and alignment weight $\lambda$, as shown in Figure~\ref{fig:hyperparameter}. The performance remains relatively stable when $\gamma$ ranges from 0.1 to 0.5, suggesting that sufficient visible patches allow the model to reconstruct masked regions by combining local variations with global information. When $\gamma$ reaches 0.6, the substantial loss of visible context disrupts temporal continuity and causes a marked performance decline. Increasing $\lambda$ from 0 to 1 improves the overall performance, demonstrating the effectiveness of explicitly modeling the patch-wise conditional distribution between semantic and temporal representations. However, increasing $\lambda$ to 2 yields no consistent improvement, suggesting that an excessive alignment weight may shift the optimization away from the reconstruction objective.

\subsection{Robustness Analysis}
\noindent We add Gaussian noise only to the test datasets to assess the robustness of \method{}, with the noise magnitude controlled by $\alpha$ and scaled by the channel-wise standard deviation. As shown in Figure~\ref{fig:test_noise_robustness}, performance decreases gradually as $\alpha$ increases. Because both numerical representations and textual semantics are derived from the same input series, noise simultaneously affects the two modalities. The relatively stable performance under moderate noise suggests that combining local variations with global information preserves sufficient context for reconstruction. Under stronger noise, both local temporal patterns and global information become increasingly distorted, reducing the reliability of constructed semantics and impairing the reconstruction of normal behavior.





\subsection{Cross-Modal Alignment Analysis}

\noindent To examine how the proposed patch-wise probabilistic cross-modal
alignment models semantic--temporal correspondences and their uncertainty
across individual temporal patches, we analyze the empirical patch-level
statistics and compare the representations learned with Gaussian NLL
alignment and contrastive alignment.
 As shown in
Figure~\ref{fig:gaussian_nll}, anomalous patches tend to exhibit higher and
more dispersed NLL values than normal patches, indicating lower likelihood
under the corresponding semantic-conditioned Gaussian distributions. The
predicted $\sigma$ varies across patches, allowing the Gaussian NLL objective
to adjust the residual penalty according to the estimated uncertainty of each
correspondence, rather than imposing an equally strong point-to-point
constraint on all patch pairs. Figure~\ref{fig:environment_similarity}
further shows that, in the visualized case, the textual and temporal
representations learned with Gaussian NLL alignment exhibit a clearer and more
consistent patch-wise correspondence pattern than those learned with
contrastive alignment.


\begin{figure}[!t]
\centering
\includegraphics[width=\linewidth]{Figures/environment_similarity_comparison.pdf}
\caption{Cross-modal cosine similarities on Environment under
(a) contrastive and (b) Gaussian NLL alignment.}
\label{fig:environment_similarity}
\end{figure}

\subsection{Detection Case Studies}

\noindent To examine how the reconstruction-based anomaly scores respond to
different anomaly manifestations, we visualize representative detection cases
from the six datasets. As shown in Figure~\ref{fig:anomaly_score_grid}, the
scores generally increase within the annotated anomalous intervals. Abrupt
deviations produce localized peaks, whereas persistent deviations lead to
continuously elevated responses, indicating that the reconstruction errors
provide distinguishable responses to these anomaly manifestations in the
selected cases.

\section{Conclusion}
\noindent We propose LoGTA, an uncertainty-aware local-to-global temporal semantic framework for multimodal time-series anomaly detection. It enriches local patch semantics with global temporal information and models semantic--temporal correspondences using  conditional Gaussian distributions, enabling uncertainty-aware patch-wise alignment. The aligned semantic representations further guide masked time-series reconstruction for anomaly detection. Experiments on six datasets demonstrate the effectiveness of LoGTA, while ablation studies validate its components.

\bibliography{aaai2027}

\end{document}
